\section{Finite-temperature Kubo transport}
\label{sec:kubo}
\subsection{Liouvillian formulation and dissipation}
Let $f(H)=[1+\exp((H-\mu)/k_BT)]^{-1}$ and define the energy-valued
Liouvillian $\mathcal L_H(A)=\ii[H,A]$. In the relaxation-time
approximation, a positive energy broadening $\Gamma$ replaces a
more general dissipative superoperator. With
$\nabla_a=-\ii[X_a,\,\cdot\,]$, a convenient conductivity convention is
\begin{equation}
 \sigma_{ab}=-\frac{e^2}{\hbar}
 \mathcal T\!\left[
   (\nabla_aH)(\Gamma+\mathcal L_H)^{-1}\nabla_b f(H)
 \right].
 \label{eq:kubo-liouville}
\end{equation}
The real symmetric part gives the dissipative DC response considered
here. The factor $e>0$ is the magnitude of the electron charge.
This expression states the operator structure behind the
finite-volume calculation, including the role of dissipation
\cite{KuehneTransport2020}. A constant $\Gamma$ is an approximation to
scattering; it is not obtained from the SCF convergence tolerance or
from a plotted density-of-states width.

Let $a^b_{nm}=\langle n|\nabla_bH|m\rangle=\hbar v^b_{nm}$, and define
the nonnegative Fermi divided difference
\begin{equation}
 F_{nm}=
 \begin{cases}
 \displaystyle-\frac{f(\varepsilon_n)-f(\varepsilon_m)}
 {\varepsilon_n-\varepsilon_m},&\varepsilon_n\ne\varepsilon_m,\\[0.7em]
 \displaystyle-f'(\varepsilon_n)
 =\frac{f(\varepsilon_n)[1-f(\varepsilon_n)]}{k_BT},
 &\varepsilon_n=\varepsilon_m .
 \end{cases}
 \label{eq:fermi-divided}
\end{equation}
A spectral expansion of Eq.~\eqref{eq:kubo-liouville} gives
\begin{equation}
 \sigma^{\mathrm{diss}}_{ab}
 =\frac{e^2}{\hbar V_d}
 \sum_{\bk,s}w_{\bk}\,g_s
 \sum_{nm}F_{nm\bk s}
 \frac{\Gamma}{\Gamma^2+(\varepsilon_{n\bk s}-\varepsilon_{m\bk s})^2}
 \Re\!\left[a^a_{nm\bk s}a^b_{mn\bk s}\right].
 \label{eq:kubo-spectral}
\end{equation}
Here $\sum_{\bk} w_{\bk}=1$, $V_d$ is the periodic length, area, or
volume, and $g_s$ accounts for spin occupancy without double counting.
For explicit SOC spinors, each state has unit maximum occupation.
The tensor is projected onto the periodic Cartesian subspace before
reporting a reduced-dimensional conductivity.

Equation~\eqref{eq:kubo-spectral} is positive semidefinite for
$T,\Gamma>0$. In particular, its diagonal elements are nonnegative.
The units are conductance times length$^{2-d}$: S\,m in one dimension,
S in two dimensions, and S/m in three dimensions, with the latter
optionally printed as S/cm. Vacuum padding must not change a
two-dimensional sheet conductivity through an arbitrary three-dimensional
volume factor. An isolated finite-system response is not automatically
a measurable bulk conductivity.

\subsection{Degenerate levels and intraband response}
The equal-energy branch of Eq.~\eqref{eq:fermi-divided} applies to all
matrix elements within a degenerate subspace, not just $n=m$.
In a periodic system, the diagonal band velocity can be nonzero and
provides an intraband term proportional to $-f'/\Gamma$.
Discarding diagonal terms would remove this contribution.
Likewise, dividing current matrix elements by
$\varepsilon_n-\varepsilon_m$ cannot be used to define a stable
position operator in a degenerate block.

A unitary change of basis inside an exactly degenerate subspace rotates
the complete current block. Because the scalar kernel is constant on
that block, the summed contraction is unchanged.
This is why comparisons should involve block traces or projectors
rather than phases of individual eigenvectors.
Near equal energies, the implementation evaluates the divided
difference with a smooth finite-temperature expansion instead of
subtracting nearly identical Fermi functions.

\subsection{Three current-operator choices}
The historical atom-embedding approximation associates each AO with
its nuclear center. It provides an economical representation but
does not contain all intra-atomic Gaussian position matrix elements.
It remains a distinct, explicitly selected approximation.

For isolated systems, the projected-AO alternative uses
\begin{equation}
 J_a=HS^{-1}X_a-X_aS^{-1}H,\qquad
 \hbar v_a=\ii J_a .
 \label{eq:projected-current}
\end{equation}
$J_a$ is anti-Hermitian. In an $S$-orthonormal eigenbasis it becomes
$(J_a)_{nm}=(\varepsilon_n-\varepsilon_m)(X_a^{\mathrm{MO}})_{nm}$.
The complex conjugation in a conductivity contraction must be retained
when SOC is present. This construction is exact for the projected
finite AO model; the difference between projected and continuum
commutators remains a basis-convergence question.

For periodic calculations, a bounded Cartesian position in one
simulation cell is not a valid replacement for the periodic current.
The Bloch path uses derivatives of the image Hamiltonian and overlap
together with analytic first-moment integrals. In atomic units its
matrix elements are
\begin{equation}
 \begin{split}
 v^a_{nm}(\bk)={}&
 [C^\dagger\partial_a H C]_{nm}
 -\varepsilon_n[C^\dagger\partial_a S C]_{nm}\\
 &+\ii(\varepsilon_n-\varepsilon_m)[C^\dagger D_aC]_{nm},
 \end{split}
 \label{eq:bloch-velocity}
\end{equation}
where $\partial_a$ differentiates Cartesian $k_a$.
The cell-gauge connection obeys
\begin{equation}
 D_a-D_a^\dagger=-\ii\,\partial_a S .
 \label{eq:connection-identity}
\end{equation}
These identities make explicit why the derivative of $H$ alone,
the canonical momentum, and the full velocity are not interchangeable
in a local-orbital calculation
\cite{EsteveParedes2023}.

Hermiticity can be checked directly: the anti-Hermitian part from
the energy-weighted overlap derivative is cancelled by
Eq.~\eqref{eq:connection-identity}.
For $n=m$, Eq.~\eqref{eq:bloch-velocity} reduces to the generalized
Hellmann--Feynman derivative
$c_n^\dagger(\partial_aH-\varepsilon_n\partial_aS)c_n$.
No eigenvector derivative or energy-denominator division is required.

For distinct band energies, the same operator determines the interband
position (dipole) matrix elements in atomic units,
\begin{equation}
 d^a_{nm}(\bk)
 =-\frac{\ii\,v^a_{nm}(\bk)}
 {\varepsilon_{n\bk}-\varepsilon_{m\bk}},
 \qquad \varepsilon_{n\bk}\ne\varepsilon_{m\bk}.
 \label{eq:interband-dipole}
\end{equation}
Here $d^a_{nm}$ has units of length and does not include the electron
charge. Substitution of Eq.~\eqref{eq:bloch-velocity} gives the
Hamiltonian-derivative term, the energy-weighted overlap-derivative term,
and $[C^\dagger D_a C]_{nm}$. All three are required in a general
nonorthogonal local-orbital basis~\cite{EsteveParedes2023}.
The k-point moment calculation constructs these matrices from real-space
image-cell integrals and can use the SCF mesh, a separate \texttt{KPOINTS}
set, or a \texttt{KPOINT\_SET} path in the moments print section.
Equation~\eqref{eq:interband-dipole} is only an interband relation.
It neither supplies diagonal positions in a periodic system nor replaces
the undivided velocity blocks needed for degenerate and intraband
transport.

\subsection{Gauge covariance and nonlocal operators}
Under a differentiable, nonsingular AO basis change $T(\bk)$,
$H'=T^\dagger HT$, $S'=T^\dagger ST$, and $C'=T^{-1}C$.
The connection must transform as
\begin{equation}
 D'_a=T^\dagger D_aT-\ii(\partial_aT)^\dagger ST .
 \label{eq:connection-transform}
\end{equation}
Substitution into Eq.~\eqref{eq:bloch-velocity} cancels all derivatives
of $T$ by using $HC=SC\varepsilon$.
The physical velocity is therefore unchanged even under a
k-dependent nonunitary AO transformation.
This is a stronger test than comparing band energies.

The Hamiltonian derivatives include all retained image contributions,
including nonlocal pseudopotential and GTH-SOC terms.
The connection term is still required; including a derivative of the
SOC Hamiltonian does not make the bare Peierls derivative complete.
SCF-potential convergence, scalar or spinor basis convergence, and
analysis-mesh convergence must be checked independently.

\subsection{Property symmetry and practical scope}
The optional Bloch-Kubo symmetry path diagonalizes only representatives
of its own analysis mesh. Atom permutations, AO rotations, lattice-image
phases, physical spin rotations, and antiunitary conjugation reconstruct
the remaining eigenframes. Each candidate frame is tested through
$C^\dagger SC=I$, $C^\dagger HC=\varepsilon$, and target/source current
covariance. An unacceptable map is not silently inserted into the
conductivity.

For a spatial Cartesian operation $R$, velocities rotate as a polar
vector. Time reversal adds the reversal of velocity as well as complex
conjugation. In the anti-Hermitian current convention used internally,
the extra factor of $\ii$ must also be conjugated.
The response sums the full tensor over the full mesh; it is not an
irreducible scalar-weight substitution for anisotropic currents.

The current implementation provides symmetric dissipative DC
conductivity with a prescribed relaxation scale. It does not yet
provide anomalous Hall, spin Hall, a general frequency-dependent
response, or microscopic scattering superoperators.
Its Bloch matrices and eigenframes are a bounded dense reference path,
although eigensystems are distributed over MPI ranks.
Fewer diagonalizations do not imply proportional memory or runtime
savings when target operators are still explicitly evaluated.
